\section{Methodology}
\label{sec:methodology}

\subsection{Problem Formulation}
Let $f_{\theta}: \mathcal{X} \rightarrow \mathcal{Y}$ denote a neural network parameterized by $\theta$, where $\mathcal{X} \subset \mathbb{R}^{d}$ is the input space and $\mathcal{Y} = \{1, \ldots, K\}$ is the label space. Given a clean input $x \in \mathcal{X}$ with label $y \in \mathcal{Y}$, an adversarial perturbation $\delta$ is crafted to maximize the loss $\mathcal{L}(f_{\theta}(x+\delta), y)$ while constrained to $\|\delta\|_{p} \leq \epsilon$.

\subsection{Adaptive Feature Alignment}
Our key contribution is the Adaptive Feature Alignment (AFA) mechanism. For a network with $L$ layers, let $h_{l}(x)$ denote the feature representation at layer $l$ for input $x$. We define the feature alignment loss as:

\begin{equation}
\mathcal{L}_{\text{align}} = \sum_{l=1}^{L} \alpha_{l} \cdot \|h_{l}(x) - h_{l}(x+\delta)\|_{2}^{2}
\end{equation}

where $\alpha_{l}$ are learnable weights that adaptively determine the importance of each layer. The weights are computed as:

\begin{equation}
\alpha_{l} = \frac{\exp(w_{l})}{\sum_{j=1}^{L} \exp(w_{j})}
\end{equation}

where $w_{l}$ are trainable parameters initialized to $1/L$.

\subsection{Training Objective}
The overall training objective combines the standard cross-entropy loss with the feature alignment term:

\begin{equation}
\mathcal{L}_{\text{total}} = \mathcal{L}_{\text{CE}}(f_{\theta}(x), y) + \lambda \cdot \mathcal{L}_{\text{align}}(x, x+\delta)
\end{equation}

where $\lambda$ is a hyperparameter controlling the strength of the alignment regularization. We use PGD-10 for generating adversarial examples during training.

\subsection{Theoretical Analysis}
We provide a theoretical justification for our approach. Consider the Lipschitz continuity of the feature mapping:

\begin{equation}
\|h_{l}(x) - h_{l}(x+\delta)\| \leq \text{Lip}(h_{l}) \cdot \|\delta\|
\end{equation}

By minimizing the feature alignment loss, we effectively constrain the Lipschitz constant of the feature mappings, leading to improved robustness.